% Part: incompleteness
% Chapter: representability-in-q
% Section: basic-representable

\documentclass[../../../include/open-logic-section]{subfiles}

\begin{document}

\olfileid{inc}{req}{bre}
\olsection{توابع پایه در~$\Th{Q}$ بازنمایی‌پذیرند}

نخست باید نشان دهیم همهٔ توابع پایه در~$\Th{Q}$ بازنمایی‌پذیرند. در
نهایت باید نشان دهیم چگونه به هر تابع پایهٔ $k$موضعی
$f(x_0,\dots,x_{k-1})$، !!a{formula}ی چون
$!A_f(x_0,\dots,x_{k-1},y)$ نسبت دهیم که آن را بازنمایی کند.

خواهیم توانست صفر، جانشین، جمع، ضرب، تابع مشخصهٔ برابری و افکنش‌ها را
بازنمایی کنیم. در هر مورد، فرمول بازنمایندهٔ مناسب کاملاً سرراست
است؛ برای نمونه، صفر با فرمول~$y=\Obj 0$، جانشین با !!{formula}
$x_0'=y$، و جمع با !!{formula}~$(x_0+x_1)=y$ بازنمایی می‌شود. کار اصلی
نشان‌دادن این است که~$\Th{Q}$ می‌تواند !!{sentence}sی مربوط را ثابت کند؛
برای نمونه، اینکه جمع با !!{formula} بالا بازنمایی می‌شود مستلزم آن است
که نشان دهیم به ازای هر زوج از اعداد طبیعی~$m$ و~$n$، $\Th{Q}$
عبارات زیر را ثابت می‌کند:
\begin{align*}
& \eq[\num n + \num m][\num {n+m}] \text{ و}\\
& \lforall[y][(\eq[(\num n + \num m)][y] \lif \eq[y][\num{n+m}])].
\end{align*}

\begin{prop}
\ollabel{prop:rep-zero}
تابع صفر~$\Zero(x)=0$ در~$\Th{Q}$ با
$!A_{\Zero}(x,y)\ident\eq[y][\Obj 0]$ بازنمایی می‌شود.
\end{prop}

\begin{prop}
\ollabel{prop:rep-succ}
تابع جانشین~$\Succ(x)=x+1$ در~$\Th{Q}$ با
$!A_{\Succ}(x,y)\ident\eq[y][x']$ بازنمایی می‌شود.
\end{prop}

\begin{prop}
\ollabel{prop:rep-proj}
تابع افکنش~$\Proj{n}{i}(x_0,\dots,x_{n-1})=x_i$ در~$\Th{Q}$ با
\[!A_{\Proj{n}{i}}(x_0, \dots, x_{n-1}, y) \ident \eq[y][x_i].\]
بازنمایی می‌شود.
\end{prop}

\begin{prob}
ثابت کنید $\eq[y][\Obj 0]$، $\eq[y][x']$ و~$\eq[y][x_i]$ به‌ترتیب
$\Zero$، $\Succ$ و~$\Proj{n}{i}$ را بازنمایی می‌کنند.
\end{prob}

\begin{prop}
\ollabel{prop:rep-id}
تابع مشخصهٔ~$=$،
\[
\Char{=}(x_0, x_1) =
\begin{cases}
  1 & \text{اگر } x_0 =x_1\\
  0 & \text{در غیر این صورت}
\end{cases}
\]
در~$\Th{Q}$ با فرمول زیر بازنمایی می‌شود:
\[
  !A_{\Char{=}}(x_0, x_1, y) \ident (\eq[x_0][x_1] \land \eq[y][\num{1}]) \lor (\eq/[x_0][x_1] \land
  \eq[y][\num{0}]).
\]
\end{prop}

اثبات به لم زیر نیاز دارد.

\begin{lem}
\ollabel{lem:q-proves-neq} اگر $n$ و~$m$ اعداد طبیعی باشند و
$n\neq m$، آنگاه $\Th{Q}\Proves\eq/[\num n][\num m]$.
\end{lem}

\begin{proof}
با استقرا روی~$n$ نشان دهید برای هر~$m$، اگر $n\neq m$، آنگاه
$Q\Proves\eq/[\num n][\num m]$.

در حالت پایه، $n=0$. اگر $m$ برابر~$0$ نباشد، برای عدد طبیعی~$k$ای
$m=k+1$. اصل موضوعه‌ای داریم که می‌گوید
$\lforall[x][\eq/[0][x']]$. با یک اصل موضوعهٔ سور، با جایگزین‌کردن
$x$ با~$\num k$، می‌توانیم $\eq/[0][\num k']$ را نتیجه بگیریم. اما
$\num k'$ همان~$\num m$ است.

در گام استقرا می‌توانیم ادعا را برای~$n$ فرض و~$n+1$ را بررسی کنیم.
بگذارید $m$ هر عدد طبیعی باشد. دو امکان هست: یا~$m=0$، یا برای
بعضی~$k$، $m=k+1$. حالت نخست مانند بالا رسیدگی می‌شود. در حالت دوم
فرض کنید $n+1\neq k+1$. آنگاه $n\neq k$. بنا بر فرض استقرا برای~$n$
داریم $\Th{Q}\Proves\eq/[\num n][\num k]$. اصل موضوعه‌ای داریم که
می‌گوید
$\lforall[x][\lforall[y][\eq[x'][y']\lif\eq[x][y]]]$. با یک اصل
موضوعهٔ سور داریم
$\eq[\num n'][\num k']\lif\eq[\num n][\num k]$. با استفاده از
منطق گزاره‌ای می‌توانیم در~$\Th{Q}$ نتیجه بگیریم
$\eq/[\num n][\num k]\lif\eq/[\num n'][\num k']$. با وضع مقدم
$\eq/[\num n'][\num k']$ را نتیجه می‌گیریم، که همان چیزی است که
می‌خواستیم، زیرا~$\num k'$ همان~$\num m$ است.
\end{proof}

\begin{explain}
توجه کنید لم ادعای چندانی ندارد: در اصل می‌گوید~$\Th{Q}$ می‌تواند
ثابت کند عددنماهای متفاوت بر اشیای متفاوت دلالت می‌کنند. برای نمونه،
$\Th{Q}$ ثابت می‌کند $0''\neq0'''$. اما نشان‌دادن این امر در حالت
کلی به کمی دقت نیاز دارد. همچنین توجه کنید با آنکه از استقرا استفاده
می‌کنیم، این استقرا \emph{بیرون} از~$\Th{Q}$ است.
\end{explain}

\begin{proof}[اثبات \olref{prop:rep-id}]
اگر $n=m$، آنگاه $\num{n}$ و~$\num{m}$ یک ترم‌اند و
$\Char{=}(n,m)=1$. اما
$\Th{Q}\Proves(\eq[\num{n}][\num{m}]\land
\eq[\num{1}][\num{1}])$؛ پس~$!A_=(\num{n},\num{m},\num{1})$ را
ثابت می‌کند. اگر $n\neq m$، آنگاه $\Char=(n,m)=0$. بنا بر
\olref{lem:q-proves-neq}،
$\Th{Q}\Proves\eq/[\num{n}][\num{m}]$ و بنابراین
$(\eq/[\num{n}][\num{m}]\land\Obj 0=\Obj 0)$ را نیز ثابت می‌کند.
پس $\Th{Q}\Proves !A_=(\num{n},\num{m},\num{0})$.

برای بخش دوم نیز دو حالت داریم. اگر $n=m$، باید نشان دهیم
$\Th{Q}\Proves\lforall[y][(!A_=(\num{n},\num{m},y)
\lif\eq[y][\num{1}])]$. به‌طور غیرصوری استدلال می‌کنیم. فرض کنید
$!A_=(\num{n},\num{m},y)$؛ یعنی
\[
(\eq[\num{n}][\num{n}] \land \eq[y][\num{1}]) \lor
(\eq/[\num{n}][\num{n}] \land \eq[y][\num{0}]).
\]
فصل نخست بنا بر منطق~$\eq[y][\num{1}]$ را نتیجه می‌دهد؛ فصل دوم با
$\eq[\num{n}][\num{n}]$ در تناقض است، و این آخری بنا بر منطق
اثبات‌پذیر است.

از سوی دیگر، فرض کنید $n\neq m$. آنگاه
$!A_=(\num{n},\num{m},y)$ عبارت است از
\[
(\eq[\num{n}][\num{m}] \land \eq[y][\num{1}]) \lor
(\eq/[\num{n}][\num{m}] \land \eq[y][\num{0}]).
\]
اینجا فصل نخست با~$\eq/[\num{n}][\num{m}]$ در تناقض است، که بنا بر
\olref{lem:q-proves-neq} در~$\Th{Q}$ اثبات‌پذیر است؛ فصل دوم
$\eq[y][\num{0}]$ را نتیجه می‌دهد.
\end{proof}

\begin{prop}
\ollabel{prop:rep-add}
تابع جمع~$\Add(x_0,x_1)=x_0+x_1$ در~$\Th{Q}$ با
\[
  !A_{\Add}(x_0, x_1, y) \ident \eq[y][(x_0 + x_1)].
\]
بازنمایی می‌شود.
\end{prop}

\begin{lem}
\ollabel{lem:q-proves-add}
$\Th{Q} \Proves \eq[(\num{n} + \num{m})][\num{n+m}]$.
\end{lem}

\begin{proof}
این را با استقرا روی~$m$ ثابت می‌کنیم. اگر $m=0$، ادعا این است که
$\Th{Q}\Proves\eq[(\num{n}+\Obj 0)][\num{n}]$. این از اصل موضوعهٔ
$!Q_4$ نتیجه می‌شود. اکنون ادعا را برای~$m$ فرض کنید؛ ادعا را برای
$m+1$ ثابت کنیم، یعنی ثابت کنیم
$\Th{Q}\Proves\eq[(\num{n}+\num{m+1})][\num{n+m+1}]$. توجه کنید
$\num{m+1}$ همان~$\num{m}'$ و~$\num{n+m+1}$ همان~$\num{n+m}'$ است.
بنا بر اصل موضوعهٔ~$!Q_5$،
$\Th{Q}\Proves\eq[(\num{n}+\num{m}')][(\num{n}+\num{m})']$. بنا بر
فرض استقرا،
$\Th{Q}\Proves\eq[(\num{n}+\num{m})][\num{n+m}]$. پس
$\Th{Q}\Proves\eq[(\num{n}+\num{m}')][\num{n+m}']$.
\end{proof}

\begin{proof}[اثبات \olref{prop:rep-add}]
!!{formula}~$!A_\Add(x_0,x_1,y)$ که~$\Add$ را بازنمایی می‌کند
$\eq[y][(x_0+x_1)]$ است. نخست نشان می‌دهیم اگر~$\Add(n,m)=k$، آنگاه
$\Th{Q}\Proves !A_\Add(\num{n},\num{m},\num{k})$؛ یعنی
$\Th{Q}\Proves\eq[\num{k}][(\num{n}+\num{m})]$. اما چون~$k=n+m$،
$\num{k}$ همان~$\num{n+m}$ است، و در~\olref{lem:q-proves-add} نشان
داده‌ایم $\Th{Q}\Proves\eq[(\num{n}+\num{m})][\num{n+m}]$.

همچنین باید نشان دهیم اگر~$\Add(n,m)=k$، آنگاه
\[
\Th{Q} \Proves \lforall[y][(!A_\Add(\num{n}, \num{m}, y) \lif
  \eq[y][\num{k}])].
\]
فرض کنید $\eq[(\num{n}+\num{m})][y]$ را داریم. چون
\[
\Th{Q} \Proves \eq[(\num{n}+\num{m})][\num{n+m}],
\]
می‌توانیم سمت چپ را با~$\num{n+m}$ جایگزین کنیم و برای~$y$ دلخواه
$\eq[\num{n+m}][y]$ را به دست آوریم.
\end{proof}

\begin{prop}
\ollabel{prop:rep-mult}
تابع ضرب~$\Mult(x_0,x_1)=x_0\cdot x_1$ در~$\Th{Q}$ با
\[
  !A_{\Mult}(x_0, x_1, y) \ident y = (x_0 \times x_1).
\]
بازنمایی می‌شود.
\end{prop}

\begin{proof} تمرین. \end{proof}

\begin{lem}
\ollabel{lem:q-proves-mult}
$\Th{Q} \Proves \eq[(\num{n} \times \num{m})][\num{n \cdot m}]$.
\end{lem}

\begin{proof} تمرین. \end{proof}

\begin{prob}
\olref[inc][req][bre]{lem:q-proves-mult} را ثابت کنید.
\end{prob}

\begin{prob}
با استفاده از~\olref[inc][req][bre]{lem:q-proves-mult}،
\olref[inc][req][bre]{prop:rep-mult} را ثابت کنید.
\end{prob}

\begin{explain}
  به یاد آورید برای نماد تابع در زبان حساب از~$\times$ و برای عمل
  معمول ضرب روی اعداد از~$\cdot$ استفاده می‌کنیم. پس~$\cdot$ می‌تواند
  میان عبارت‌های مربوط به اعداد ظاهر شود (مانند~$m\cdot n$)، درحالی‌که
  $\times$ فقط میان ترم‌های زبان حساب می‌آید (مانند
  $(\num{m}\times\num{n})$). گیج‌کننده‌تر اینکه~$+$ هم برای !!{function}
  و هم برای عمل جمع به کار می‌رود. هنگامی که میان ترم‌ها ظاهر شود---
  برای نمونه در~$(\num{n}+\num{m})$---!!{function} دوموضعیِ زبان حساب
  است، و هنگامی که میان اعداد بیاید---برای نمونه در~$n+m$---عمل جمع
  است. این شامل حالت~$\num{n+m}$ نیز می‌شود: این عددنمای استانداردِ
  متناظر با عدد~$n+m$ است.
\end{explain}

\end{document}
